\chapter{Counterfoil Choice: A Cultural Technology for Training Admissibility Recognition}
\label{ch:counterfoil}

\begin{plainlanguage}
\textbf{In plain terms:} Among the Kalabari people, a person is sometimes offered two options where only one is ever meant to be chosen — the other exists purely to demonstrate what failure looks like. This chapter argues the device is a training contrast: making the wrong answer absurdly wrong lets a learner cheaply practice recognizing the \emph{shape} of inadmissibility before facing subtler, higher-stakes cases.
\end{plainlanguage}

\begin{chapternote}
This part exists to do something the previous one could not: find a domain where the elimination machinery of Part~II is not merely tested against messy data, but where something resembling it is already, independently, a cultural technology --- practiced, named, and reflected upon by the people who use it, centuries before this book's formalism existed to describe it.
\end{chapternote}

\section{A tool the Kalabari already built}

In 1999, Nimi Wariboko gave a name to something Kalabari elders had been doing for generations without needing to formalize it. He called it \emph{counterfoil choice}: a structure in which two options are presented, but only one is a genuine alternative. The other --- the counterfoil --- is offered specifically to be refused. A father places a full bowl of garri and an empty bowl before his child. A bride is served an elaborate wedding feast and ritually turns her face from it. A young soldier is asked to throw his lame general at the enemy like a human cannonball. In each case, Wariboko writes, $A$ is not an alternative to $B$; $B$ is only a counterfoil to $A$. When an option is offered, the negative is put forward to show the person where the decision should not go. The decision-maker retains the formal freedom to choose either option, but the outcome is supposed to be foreclosed in advance --- the freedom to choose is not limited the way it is in a Hobson's choice, yet the result of the selection resembles a Hobson's choice all the same.

Wariboko is explicit that this device is not, in his framing, a test of intelligence. He calls it a \emph{culture wisdom quotient} test, designed to probe a person's internalization of what he terms the Kalabari aristocratic ideal --- stylishness, nonchalant achievement, self-knowledge (\emph{bu nimi}) --- not their capacity to think. That disclaimer is the most useful sentence in the source material, because it draws a boundary this chapter needs and will respect rather than override: the question is not what the Kalabari say the practice tests. It is what the practice, examined as a piece of pedagogical engineering, actually trains --- and the answer, this chapter will argue, is something prior to and more general than either intelligence or wisdom as ordinarily understood.

\section{The counterfoil as a high-margin training contrast}

A child does not learn the concept ``edible'' by being shown a single bowl of food and told its name. Wariboko's examples instead pair the admissible item with a deliberately, almost comically inadmissible one: a bucket against a basket that cannot hold water; a full bowl against an empty one. The distance between the admissible option and the counterfoil is not incidental. It is the entire pedagogical point --- the lesson works precisely because that distance is made as large as possible.

\begin{definition}[Counterfoil Pair]
A counterfoil pair is a tuple $(a,c)$ where $a$ is the admissible option, $c$ is the inadmissible option, and the pair is selected so as to maximize $d(a,c)$, the distance between them, subject to remaining pedagogically relevant to the lesson at hand.
\end{definition}

What is being installed by repeated exposure to such pairs is not the content of any single lesson --- children do not need extensive practice to know a basket cannot hold water once the affordance is understood. What is being installed is something more general: the \emph{shape} of a non-option. A child who has been shown enough wildly-mismatched pairs across enough domains --- vessels, food, conduct --- begins to acquire a transferable competence: the ability to recognize, quickly and with low cognitive cost, that something fails to satisfy a goal. The counterfoil is training wheels for a classifier, not a lesson about any particular bowl.

This reframes the formal object of inquiry. The interesting structure here is not a shrinking space of candidate definitions being narrowed by elimination, in the sense Part~II built --- that mapping was tried earlier in this project's development and abandoned as overreaching once it became clear the Kalabari material was not actually modeling definitional convergence. The interesting structure is a developing function
\[
f : X \to \{0,1\}
\]
mapping candidate options to admissible or inadmissible, trained initially on pairs $(a,c)$ chosen so that $d(a,c) \gg 0$, on the implicit pedagogical bet that competence built on exaggerated contrasts will eventually transfer to judgments made under much smaller margins, where the boundary is no longer obvious and must instead be inferred.

\begin{analogy}[demoted from ``Margin Theorem'']
Counterfoil pedagogy resembles high-margin training in statistical classification, where larger separation between classes is associated empirically with faster or more robust acquisition of a decision boundary.
\end{analogy}

This is stated explicitly as an analogy and not a theorem: the formal guarantees underlying real margin-based learning results --- assumptions about convexity, VC-dimension bounds, independent and identically distributed sampling --- have no established analogue for how a human child acquires categories through cultural pedagogy. A genuine theorem here would require either a formal model of human concept acquisition under exaggerated contrast, which does not currently exist, or direct empirical evidence that children's learning rate scales with the size of a training contrast, which has not been cited in support of this chapter. The analogy is retained because it is evocative and organizes the intuition usefully --- not because it has been proven to hold.
